% Paper A standalone outline prototype
% Version: 2026-07-13
% Status: standalone LaTeX outline prototype; local Overleaf has an eight-chapter input skeleton, but full prose/results migration is not claimed.
% Boundary: does not rewrite historical preregistration, protocol, assignment,
% routing, analysis artifacts, or the current Overleaf body.
% The actual manuscript source is the local main.tex plus sections/*.tex.
% This file is not directly input-able into that manuscript project.
\documentclass[UTF8,a4paper,12pt]{ctexart}
\usepackage[a4paper,margin=2.4cm]{geometry}
\usepackage{amsmath,amssymb,booktabs,longtable,array,enumitem,hyperref}
\hypersetup{hidelinks}
\setlist{nosep}
\newcommand{\contract}[5]{\paragraph{写作合同。}问题：#1；证据：#2；primary：#3；sensitivity/audit：#4；禁止主张：#5。}
\newcommand{\primary}{\textbf{primary}}
\newcommand{\sensitivity}{\textbf{sensitivity/audit}}

\title{Paper A 正式论文提纲：\texorpdfstring{$R_u$}{Ru}--\texorpdfstring{$D_u$}{Du} 双链路与冻结路由}
\author{写作合同 v3}
\date{2026-07-13}

\begin{document}
\maketitle

\begin{quote}
本文件是 Paper A 的独立 LaTeX 结构原型；Markdown 版本是章节与主张审计合同。实际正文真源是本地 manuscript 的 `main.tex + sections/*.tex`，本文件不能代替或直接输入正文。正式主线固定为
\texttt{Pilot -> PreScreen -> Calibration -> Main(Test + Validation)}，轮次固定为
\texttt{P1 -> C1 -> C2 -> T1 -> V1}。本文件不回写历史预注册、protocol freeze、assignment、routing、统计执行参数或任何原始工件；当前 Overleaf 仅建立八章输入骨架并完成限域修正，不声称正文、结果或 PDF 已全面迁移。
\end{quote}

展示层符号固定为：$R_u$ 为 Calibration-only protocol reliability；$D_u=(G_u,S_u,C_u,V_u,P_u)$ 为 P1-informed diagnostic profile，五维均 higher-is-better；$B_u/F_u/Q_u$ 分别为 blind-trust、condition-fragility、computability/geometry-gate risk，均 lower-is-better 且不属于 $D_u$；$T_u/U_u$ 仅作 raw-risk legacy aliases。

\section*{核心主轴与贡献边界}

\[
\text{协议证据}\to\text{evidence validity gate}\to
\{R_u^{\mathrm{calib}},D_u\}\to\text{C2 frozen worker state}
\to\text{predictive validity}\to\text{Random/Global/Full}.
\]

一级贡献仅为：
\begin{enumerate}
  \item 阶段化、轮次化、可审计的半自动全景布局标注协议；
  \item 双链路 worker modeling：Calibration-only protocol reliability $R_u$ 与 P1-informed 多维诊断画像 $D_u$；
  \item Validation 前冻结 worker state 后的预算感知路由评价。
\end{enumerate}

support-aware 场景可靠度与 Global fallback、加权共识消融、failure-family/counterexample bank、predictive validity、active-time integrity audit 为二级创新。模型重训练、Bi-Layout、A-line Manhattan、自动 OOS、worker-facing Manhattan correction 和数据集重标注模型只能进入讨论、附录或未来工作。

\section{引言}
\subsection{问题背景与现实约束}
\contract{全景布局标注为何需要同时处理效率、质量和审计性}{任务成本、模型初始化、worker 行为与 provenance}{协议问题定义}{工程背景与收集限制}{模型重训练或 Manhattan 是本文问题核心}
\subsection{研究缺口与本文路线}
\contract{现有方法为何不足以支持冻结前可审计的 worker state 与预算路由}{相关工作对比与本文证据链}{双链路和冻结路线}{历史方案与失败案例}{正文已完成迁移}
\subsection{研究问题与贡献层级}
\contract{三个一级 RQ 及 RQ3a/RQ3b 子问题如何对应后文}{RQ 合同与贡献清单}{三项一级贡献}{二级创新}{把二级创新写成一级贡献}
\subsection{研究范围与禁止主张}
\contract{Paper A 明确不研究什么}{协议边界与支线材料}{范围声明}{附录/未来工作}{把支线写成主实验}

\section{相关工作}
\subsection{全景布局估计与布局标注}
\contract{本文任务与布局标注文献的关系}{任务定义与文献}{Paper A 定位}{指标兼容性}{A-line 工具是本文主方法}
\subsection{模型辅助标注与 automation bias}
\contract{blind trust、issue recognition 与 correction 为什么分开}{相关 bias 与 assisted annotation 文献}{RQ2 概念基础}{model\_issue 标签审计}{选择 model\_issue 即完成几何纠错}
\subsection{众包可靠度与 worker modeling}
\contract{为什么拆分 protocol reliability 与 diagnostic profile}{可靠度、画像和 consensus 文献}{双链路定位}{weighted consensus}{P1 或 semi 直接生成 $R_u$}
\subsection{自适应冗余、路由与预算评价}
\contract{为什么使用固定预算 offline replay}{routing 与 sequential redundancy 文献}{RQ3b 方法学基础}{V1 shadow/deployment}{V1 单独承担主比较}
\subsection{Provenance、完整性与可审计协议}
\contract{evidence validity 为何是测量前提}{provenance 与审计文献}{gate 合同}{post-closeout audit}{missing evidence 是 success}

\section{研究协议与数据生命周期}
\subsection{阶段、轮次与冻结点}
\contract{P1/C1/C2/T1/V1 各自允许什么}{protocol 与 assignment SOP}{主线、职责、freeze points}{extension/replication cohort}{改变 protocol freeze}
\subsection{任务池、条件、worker 与数据流}
\contract{pool/condition/assignment 如何隔离}{manifest、task pool、provenance}{阶段--数据--用途矩阵}{降级路径}{改变 assignment 或 routing}
\subsection{Scope、reference 与 final-gold provenance}
\contract{哪些任务可进入 hard geometry 或 scope evidence}{scope、final-gold、reference path/SHA}{reference gate}{soft ambiguous/unavailable}{把单一 reference 强行当唯一真值}
\subsection{Active-time identity 与时间降级}
\contract{哪些时间证据可进入 RQ1 primary}{worker--task--annotation 的 owner-valid exact browser logs、independence、合法 script/version}{exact active time}{known-only、task-level、lead\_time、unknown、parent-derived；session 内 max、跨 session 合法分段累加}{active-time 是 worker quality；跨 worker 取最大 task time}
\subsection{Evidence validity and post-closeout integrity gate}
\contract{每行 evidence 可进入哪条链路}{independence、timing、scope/final-gold、reference、process/system、expert adjudication、artifact path/SHA/rule version}{纳入 $R_u$ 或 $D_u$ 的资格}{pending、system issue、forensic rows}{missing evidence 补成 success}
\subsection{Post-closeout 修正与不回写边界}
\contract{P1 closeout 后 integrity evidence 如何使用}{只读 correction/provenance layer}{诊断与 predictive validity}{confirmed/suspected non-independence}{回写 admission、assignment、routing 或原始工件}

\section{测量模型与双链路工人画像}
\subsection{指标层级与方向}
\contract{工作量、质量、可靠度和风险率分别测什么}{指标字典与 inclusion gate}{每一 estimand 的方向、分母、support、freeze stage}{不可合并指标的审计}{把 failure rate 当 reliability}
\subsection{Reference-gated post-closeout geometry metric}
\contract{何时可报告诊断性几何 metric}{hard-single、hard-multi、pairing、范围、奇数点、歧义和 cardinality gate}{兼容组件的诊断几何}{soft/不兼容组件}{geometry score 是 GT correctness 唯一替代；A-line Manhattan 进主实验}
\subsection{链路 A：Calibration-only protocol reliability $R_u$}
\contract{正式 Calibration protocol 中 worker 的可靠度是多少}{仅 C1/C2 Calibration\_manual}{$R_u$、LCB、CI、support、freeze}{候选 $R_{u,s}$ 与 precision audit}{P1、Calibration\_semi、C2b、Main 回流}
\subsection{场景特异 $R_{u,s}$ 与 Global fallback}
\contract{何时启用场景特异 reliability}{C1/C2 scene support 与 CI}{support 达标的 $R_{u,s}$}{activation/degeneration}{support 不足时静默启用}
\subsection{链路 B：多维诊断画像 $D_u$}
\contract{worker 的可解释行为维度是什么}{经过 gate 的 P1-informed 与独立后续 evidence}{越高越好的五维画像及可执行 $G_u$ component}{component/support audit、$B_u/F_u/Q_u$ raw risk}{失败率与 reliability 共用符号}

\[
D_u=(G_u,S_u,C_u,V_u,P_u),\quad
C_u=1-\text{semi correction failure rate},\quad
V_u=1-\text{undercoverage failure rate},\quad
P_u=1-\text{process failure rate}.
\]
\[
B_u=\text{blind-trust risk},\qquad F_u=\text{condition fragility},\qquad Q_u=\text{computability/geometry-gate risk},
\]
三者均为 lower-is-better risk signatures，不进入 $D_u$；旧 $T_u/U_u$ 只保留为 legacy aliases。

\subsection{Issue recognition 与 geometry correction}
\contract{识别问题是否等于完成几何纠错}{model\_issue 与最终 geometry 的独立对照}{有 geometry evidence 时的 $C_u$}{仅 issue-recognition rows}{把 model\_issue 选择写成 correction 完成}
\subsection{Failure-family 与 counterexample bank}
\contract{如何解释失败而不污染 $R_u$}{worker-task evidence 与 expert-reviewed candidates}{五个一级 family 长表}{subfamily、insufficient cell、counterexample audit}{failure-family 等同 $R_u$ 或作为唯一贡献}
\subsection{Predictive-validity gate}
\contract{P1-informed profile 是否与独立后续行为呈描述性预测关联}{三组 predictor--target：P1 geometry→C1/C2 geometry、P1 semi-correction→Calibration\_semi correction、P1 process warning→later process reliability}{有效 worker、support、Spearman/Kendall、worker bootstrap 95\% CI、directional consistency}{suspected/not-evaluable、discrepancy workers、missing、interpretation\_allowed}{P1 自证、directional check 冒充正式 RQ3a 或把 predictive validity 当 routing utility}

\section{路由策略与统计分析}
\subsection{冻结 worker state 与 task risk}
\contract{C2 后哪些 state 固定}{C2 snapshot、scene contract、risk manifest}{worker state、support、fallback、stop rule}{activation/degeneration}{V1 改 tier、Score、$\tau_d$、$k_0/k_{\max}$ 或 stop rule}
\subsection{Random、Calibration-only Global 与 Full}
\contract{三种 policy 如何在同一预算下比较}{同一 candidate pool、预算和 sequential rule}{三策略主比较}{diagnostic-feature-assisted Full 的 calibration-only sensitivity}{未过 gate 的 P1 profile 直接 routing}
\subsection{Calibration\_manual offline replay 与 V1 隔离}
\contract{如何避免真实标签造成泄漏}{Calibration\_manual image-level cross-fitting、shadow support 与 V1 frozen deployment logs}{offline 主比较}{V1 deployment/audit}{V1 临时并跑未注册策略或把 replay 写成无偏在线效应}
RQ3b 的 cross-fitting 按 \texttt{base\_task\_id/image\_id} 划分 folds；同一 base image 的全部 annotation 同 fold。评价 fold 不参与 worker state、Score、threshold、normalization、support gate、fallback 或 stopping-rule estimation，feature selection 只在训练 folds 完成；评价只在观测候选集合 $U_t$ 内进行并报告 policy support/unsupported。Operational frozen state 供 V1，cross-fitted shadow state 仅供 replay。
\subsection{RQ1：效率统计}
\contract{半自动是否降低 owner-valid exact active time}{T1 paired condition 与 exact logs}{已冻结 RQ1 estimand、CI、配对/cluster-aware inference}{fallback、lead\_time、coverage audit}{把 lead\_time 混入 primary 或把快解释为可靠}
\subsection{RQ2：质量、纠错与行为统计}
\contract{半自动如何影响 quality、correction 和 failures}{T1/C1/C2 quality、reference、failure-family}{已冻结 quality/IAA inference}{weighted consensus、failure distribution、counterexamples}{weighted consensus 成为主轴}
\subsection{RQ3a：画像效度统计}
\contract{P1 profile 是否与后续独立行为呈描述性预测关联}{三组 gated P1→C1/C2/Main rows}{小样本 Spearman/Kendall 与 worker bootstrap CI}{support/pending/discrepancy audit；不进入 Bonferroni 六个主检验}{P1 自证}
\subsection{RQ3b：路由效用统计}
\contract{Full 是否优于 Random/Global}{Calibration\_manual image-level cross-fitted replay；fold 按 base\_task\_id/image\_id 且同图 annotation 同 fold}{RQ3b-1 首选 worker 条件性质量与 RQ3b-2 序贯预算/最终 medoid 或 LOO-compatible consensus}{policy support/unsupported、V1 deployment、fallback、activation；无独立 reference 时最终聚合仅 descriptive}{V1 作为主因果比较场；同一 submission 同时作输出和 reference}
\subsection{Missingness、sensitivity 与 audit}
\contract{缺失证据如何影响解释}{missingness、fallback、system issue、not\_evaluable}{主估计保持冻结}{降级与审计}{静默过滤或把 NA 当 success}

\section{实验结果}
\subsection{Evidence completeness 与 validity-gate 流量}
\contract{有多少证据真正可用}{各阶段/池 evidence count、status、path/SHA/rule version}{纳入主分析的 rows}{pending、system issue、not\_evaluable}{未过 gate 计为 success}
\subsection{RQ1：效率结果}
\contract{效率差异是否成立且质量未下降}{exact active-time、coverage、condition、quality-adjusted checks}{primary exact estimate}{fallback 与 fast-low-quality/blind-trust audit}{更快等于更可靠}
\subsection{RQ2：质量、纠错与 failure-family 结果}
\contract{质量和行为是否变化}{geometry、IAAt、issue recognition、correction、blind trust、undercoverage、family}{主质量/行为结果}{weighted consensus 与 counterexample}{失败 family 取代 $R_u$}
\subsection{RQ3a：跨阶段画像效度结果}
\contract{画像能否与独立后续行为呈描述性预测关联}{gated P1 与 C1/C2/Main}{三组 association、CI、directional consistency}{support/pending/discrepancy}{P1 自证}
\subsection{RQ3b：Random/Global/Full 路由结果}
\contract{固定预算下 Full 是否有用}{Calibration replay 与 V1 frozen deployment}{offline 主比较}{V1 audit/fallback/activation}{V1 重新调参}
\subsection{二级创新与审计结果}
\contract{辅助机制如何解释主结果}{scene fallback、weighted consensus、failure/counterexample、active-time audit}{补充解释}{sensitivity/audit}{把辅助结果升级一级贡献}

\section{讨论与局限}
\subsection{主要发现与协议意义}
\contract{哪些结论由 primary evidence 支持}{第6章 primary 结果}{主结论}{审计与降级}{超出 support 做强主张}
\subsection{双链路边界与 predictive validity}
\contract{$R_u$ 与 $D_u$ 如何互补且不可互换}{第4章 gate 与第6章 predictive 结果}{概念边界}{pending/insufficient}{把画像当正式 reliability}
\subsection{路由效用与部署解释}
\contract{offline replay 与 V1 各自能说明什么}{Calibration replay、V1 deployment}{RQ3b 解释}{shadow/fallback}{V1 反推 C2}
\subsection{支持度、缺失与外部效度}
\contract{support 和 missingness 如何限制结论}{support、not\_evaluable、fallback}{局限边界}{sensitivity/audit}{缺失等于成功}
\subsection{支线内容与未来工作}
\contract{模型重训练、Bi-Layout、Manhattan、自动 OOS 和 worker-facing correction 放在哪里}{支线材料与工程审计}{未来工作}{附录说明}{将其写成 Paper A 一级贡献}

\section{结论}
\subsection{三个一级 RQ 及 RQ3a/RQ3b 子问题的证据边界}
\contract{逐项回答 RQ1、RQ2、RQ3（分别报告 RQ3a 与 RQ3b）}{第6章 primary/sensitivity/audit}{不超证据的结论}{降级规则}{将 audit 当 primary}
\subsection{一级贡献与不主张事项}
\contract{本文最终贡献是什么}{协议、双链路、冻结路由}{三项一级贡献}{二级创新}{支线升级}

\appendix
\section{附录合同}
保留 artifact field contract、完整指标字典、阶段--用途矩阵、统计细节、failure subfamily、场景启用门槛、counterexample provenance、文件版本与复现入口。A-line Manhattan、模型重训练和其他支线只能作为明确标记的附录/未来工作内容。

\section*{正式图表清单}
\begin{itemize}
  \item 图1：\texttt{Pilot -> P1 -> C1 -> C2 -> T1 -> V1} 与 freeze points。
  \item 图2：evidence source $\to$ validity gate $\to R_u/D_u \to$ frozen worker state $\to$ routing。
  \item 图3：Random/Global/Full、Calibration replay、T1、V1 数据隔离。
  \item 表1：阶段--数据--允许用途矩阵。
  \item 表2：指标字典（名称、符号、方向、分子/分母、来源、gate、support、primary/sensitivity/audit、freeze stage）。
  \item 表3：RQ--数据--估计量--统计方法--降级规则。
  \item 表4：worker profile main matrix、$R_u$、$G_u/S_u/C_u/V_u/P_u$、$B_u/F_u/Q_u$、raw rates、support、legacy aliases 与 freeze stage。
  \item 表5：failure-family/subfamily 长表、support、\texttt{interpretation\_allowed} 与 expert-review status。
\end{itemize}

\end{document}
